% !PRINTED 576
\weberpage{576}

Let $\mu$ and $\nu$ be two positive integers and $x$ a
variable. We define a rational function of $x$ by
the equation

\weq{4}{(\mu, \nu) = \frac{(1 - x^{\mu})(1 - x^{\mu - 1}) \ldots (1 - x^{\mu - \nu + 1})}{(1 - x)(1 - x^2) \ldots (1 - x^{\nu})}.}

If $\nu = \mu + 1$, this function vanishes
identically, on account of the factor $1 - x^{\mu - \nu + 1}$, and likewise when $\nu$
becomes still greater.

We set up for these functions some recursion formulae,
which follow immediately from an inspection of formula (4):

\weq{5}{(1 - x^{\mu - \nu + 1})(\mu + 1, \nu) = (1 - x^{\mu + 1})(\mu, \nu)}

\weq{6}{(1 - x^{\nu + 1})\ (\mu, \nu + 1) = (1 - x^{\mu - \nu})(\mu, \nu),}

and by the combination of (5) and (6), after $\mu$ has been
replaced by $\mu + 1$ in (6),

\weq{7}{(1 - x^{\nu + 1})(\mu + 1, \nu + 1) = (1 - x^{\mu + 1})(\mu, \nu).}

If one subtracts (6) from (7) and divides by the non-vanishing $1 - x^{\nu + 1}$, it follows further

\weq{8}{(\mu + 1, \nu + 1) = (\mu, \nu + 1) + x^{\mu - \nu}(\mu, \nu).}

If we set, by way of definition, $(\mu, 0) = 1$, and $(\mu, \nu) = 0$ for
negative $\nu$, then these formulae all hold generally for every
integer $\nu$.

From the last formula we conclude by complete
induction that $(\mu, \nu)$ is an \gesperrt{integral function} of $x$. For
if this property is proved for any one $\mu$, it follows for
the next greater $\mu$ from formula (8). But, for example,
$(1, \nu)$ is equal to 1 or equal to 0, and so the theorem is generally correct.

If we put $\mu + 1$ in place of $\mu$ in (8), we can
with its help transform formula (5) thus,

\weq{9}{(1 - x^{\mu + 1})(\mu, \nu) = (\mu + 1, \nu) + (\mu + 1, \nu + 1) - (\mu + 2, \nu + 1).}

We now consider an integral function

\weq{10}{f(x, \mu) = 1 - (\mu, 1) + (\mu, 2) - (\mu, 3) + \cdots = \overset{\nu}{\Sigma} (-1)^{\nu} (\mu, \nu),}

determined by \gesperrt{Dirichlet}, \gesperrt{Cauchy}, \gesperrt{Kronecker}. The path that
\gesperrt{Gauss} takes, which we follow in the text, uses only algebraic
means. \gesperrt{Kronecker} has given in Liouville's Journal Ser.\,2, Bd.\,1 a
way of determining the sign which, with the use of a
remark of \gesperrt{Dedekind} (Schlömilch's Ztschr. Bd.\,15, Literaturzeitung
S.\,21), likewise becomes a purely algebraic one.
%CONT>>
